% TOP_PROBLEM: {"problemId": "problem.sarnak-s-mobius-randomness-conjecture-for-zero-entropy-systems", "canonicalTitle": "Sarnak's Mobius randomness conjecture for zero-entropy systems", "releaseRank": 84, "primaryDomain": "probability_ergodic_dynamics", "primaryDomainLabel": "Probability, ergodic theory & dynamics", "displayStatus": "open", "releaseStatus": "open"}
% !TeX program = lualatex
% Definition and sources only; this file is not an LLM solution attempt.
\documentclass[11pt]{article}
\usepackage[margin=25mm]{geometry}
\usepackage{fontspec}
\setmainfont{Latin Modern Roman}
\usepackage{amsmath,amssymb,mathtools}
\usepackage{unicode-math}
\setmathfont{Latin Modern Math}
\usepackage{xurl,hyperref}
\hypersetup{colorlinks=true,urlcolor=blue}
\setcounter{secnumdepth}{0}
\setlength{\parindent}{0pt}
\setlength{\parskip}{5pt}
\setlength{\emergencystretch}{3em}
\begin{document}
\section*{\texorpdfstring{Sarnak's Mobius randomness conjecture for zero-entropy systems}{Sarnak's Mobius randomness conjecture for zero-entropy systems}}\label{84}
\subsection{Definitions and mathematical statement}
The Möbius function is $\mu(n)=0$ when a prime square divides $n$, and $\mu(n)=(-1)^r$ when $n$ is a product of $r$ distinct primes. A compact metric system consists of a continuous map $T:X\to X$; zero topological entropy means its orbit complexity has zero exponential growth rate. The assertion is
\[
 \frac1N\sum_{n=1}^{N}\mu(n)f(T^nx)\longrightarrow0
 \qquad(f\in C(X),\ x\in X).
\]
Every point is quantified, not just almost every point under a chosen invariant probability measure. The test functions may be complex-valued.
\par\smallskip\textit{Formulation and concept references:} \href{https://arxiv.org/abs/2603.11087}{[S1]}.

\subsection{Short English statement}
Is the Möbius function asymptotically uncorrelated with every continuous observable of every zero-entropy dynamical system?
\subsection{Sources}
\begin{itemize}
\item[S1]Q. Liu, J. Ma, and H. Wang, The Mobius Disjointness Conjecture on infinite-dimensional torus, Introduction (2026).
\url{https://arxiv.org/abs/2603.11087}.
\textit{Formulation source.}
\item[S2]Minimal zero entropy subshifts can be unrestricted along any sparse set.
\url{https://www.cambridge.org/core/journals/ergodic-theory-and-dynamical-systems/article/minimal-zero-entropy-subshifts-can-be-unrestricted-along-any-sparse-set/1357B8C6E5361EC6439CEA52634BD7A0}.
\textit{Further reference; not a proof endorsement.}
\end{itemize}
\end{document}
